\documentclass[11pt]{article}
\usepackage[utf8]{inputenc}
\usepackage[T1]{fontenc}
\usepackage{geometry}
\usepackage{hyperref}
\usepackage{longtable}
\usepackage{array}
\usepackage{enumitem}
\usepackage{listings}
\geometry{margin=1in}

\title{Cultural Ontology for Utterance Lowering:\\
A Qualitative Meta-Analysis for Weksa}
\author{Weksa Research Notes}
\date{2026-06-09}

\begin{document}
\maketitle

\begin{abstract}
This paper synthesizes contemporary work on linguistic relativity,
sociolinguistic indexicality, language ideology, dialect grammar, and automated
language-system bias in order to define a typed cultural-ontology state shape
for Weksa. The synthesis rejects strong linguistic determinism and instead
adopts a weaker but operationally useful position: language varieties and
culture-bearing speech communities do not determine thought, but they do make
some distinctions grammatically obligatory, pragmatically salient, socially
indexed, or interactionally dangerous. For a generation system, those pressures
must be represented explicitly before interlingua is lowered into plausible
utterances. The proposed schema treats cultural ontology as a realization lens:
it does not own source meaning, speaker identity, or final wording; it owns
salience axes, social-indexical meanings, dialect affordances, activation gates,
evidence, and forbidden flattenings.
\end{abstract}

\section{Status and Scope}

This is a qualitative meta-analysis rather than a numerical meta-analysis. The
source literature spans cognitive linguistics, linguistic anthropology,
sociolinguistics, dialectology, corpus documentation, and AI fairness research.
The evidence is therefore not commensurable as effect sizes. The useful output
for Weksa is a design synthesis: a schema that can preserve cultural and
dialectal pressure without turning culture into stereotypes or dialects into
surface accent filters.

The immediate design targets are:

\begin{itemize}[noitemsep]
  \item African American English / African American Language as a dialect or
  speech tradition with systematic grammar and cultural indexicality.
  \item Regional and contact-influenced US English targets, including West
  Coast, East Coast, Deep South, and New York/Yiddish-contact registers as
  explicit overlays rather than one generic American flavor.
  \item Scottish targets that distinguish Scots, Scottish Standard English,
  Highland English, and Scottish Gaelic influence instead of collapsing them
  into a single ``Scottish'' costume.
  \item Regional Brazilian Portuguese targets such as urban Rio slang, a Sao
  Paulo television-reporter register, and a Minas Gerais worker persona, all as
  explicit profile stacks rather than one generic Brazilian flavor.
  \item Japanese regional and historical targets, including contemporary
  dialects and historical stages represented through NINJAL and related corpus
  resources.
  \item A general cultural-ontology profile format suitable for CultCache
  storage, Qdrant evidence retrieval, and later utterance lowering.
\end{itemize}

\section{Method}

Sources were selected by relevance to four questions:

\begin{enumerate}[noitemsep]
  \item What remains useful in contemporary interpretations of Sapir--Whorf or
  linguistic relativity?
  \item How do dialects, registers, and styles carry social meaning?
  \item What must be represented to produce AAVE/AAL and Scottish dialectal
  output without caricature?
  \item What failures do automated language systems already exhibit around
  dialects, especially African American English?
\end{enumerate}

The evidence was grouped into five families: relativity, indexicality and
identity, language ideology, dialect documentation, and computational bias.
Open PDFs were cached under
\texttt{docs/research/cultural-ontology/sources/pdfs/} when the source exposed a
stable, legally accessible copy. Sources that could not be cached are recorded
in the source manifest.

\section{Findings}

\subsection{Relativity Is Best Treated as Realization Pressure}

Strong linguistic determinism is not a sound engineering doctrine. Contemporary
work is more modest and more useful: language influences cognition most
reliably where speakers must package thought for expression. Slobin's
``thinking for speaking'' reframes relativity as the routinized attention that
speakers bring to events while preparing them for grammar and discourse
\cite{slobin1996}. Lucy's reformulation also separates linguistic diversity
from determinism and treats relativity as an empirical question about how
language-specific categories guide interpretation \cite{lucy1992,lucy1997}.
Wolff and Holmes summarize the modern resurgence as contested but productive:
language can affect thought, but effects vary by domain, task, and mechanism
\cite{wolffholmes2011}.

For Weksa, this means the interlingua should continue to own source meaning,
while cultural ontology owns the distinctions that become salient during
utterance realization. A cultural profile should not claim that a speaker
cannot think outside its categories. It should claim that a lowerer must inspect
some categories before producing an honest utterance.

\subsection{Culture Reaches Language Through Indexicality, Ideology, and Practice}

Language forms do not merely encode propositions. They index social relations,
stance, identity, authority, authenticity, and membership. Silverstein's
indexical order explains how forms can presuppose and create social context at
multiple orders of awareness \cite{silverstein2003}. Eckert's third-wave
variationism treats variation as a resource for social meaning, not merely a
statistical residue of demographic categories \cite{eckert2012}. Bucholtz and
Hall frame identity as interactionally produced through linguistic and other
semiotic practices, including labels, stances, styles, and structures
\cite{bucholtzhall2005}.

Language ideology adds a second warning. Irvine and Gal's mechanisms of
iconization, fractal recursivity, and erasure show how communities and observers
turn linguistic differences into social facts, often by hiding internal
variation or projecting ideology back onto language \cite{irvinegal2000}.
Therefore a Weksa schema must distinguish:

\begin{itemize}[noitemsep]
  \item a form's grammatical function;
  \item its indexical meaning in a speech community;
  \item the ideology outsiders attach to it;
  \item whether the current speaker has standing to use it.
\end{itemize}

\subsection{AAVE/AAL Must Be Modeled as Grammar Plus Social Authority}

African American English, also discussed as African American Language, is a
systematic variety with its own grammatical, phonological, lexical, and
discourse patterns. Rickford's work treats AAVE as structurally rich and
educationally significant \cite{rickford1999,rickford2016}. Green's linguistic
introduction provides a grammatical account of African American English
\cite{green2002}. The Yale Grammatical Diversity Project documents features
such as invariant or habitual \emph{be}, while ORAAL and CORAAL provide
linguistic-pattern and corpus anchors for real data rather than stereotype
\cite{yaleInvariantBe,oraalPatterns,coraal}.

The schema implication is blunt: AAVE/AAL cannot be represented as slang,
misspelling, or accent. It needs a dialect profile with:

\begin{itemize}[noitemsep]
  \item tense--modality--aspect affordances such as habitual \emph{be},
  stressed \emph{BIN}, completive \emph{done}, and negative concord;
  \item discourse and stance patterns grounded in source evidence;
  \item regional and speaker variation;
  \item activation gates based on speaker identity, scene, community, and
  authorial permission;
  \item anti-caricature constraints.
\end{itemize}

\subsection{US English Is a Regional, Social, and Contact-Language Field}

``US English'' is not a single cultural ontology. It is a base locale with
regional dialects, ethnic and racialized speech traditions, institutional
registers, class-indexed styles, media norms, and contact-language repertoires.
The Dictionary of American Regional English documents regional words, phrases,
and pronunciations across the United States, while the Atlas of North American
English provides a phonological map of North American regional sound change
\cite{dare,dareReport,anae}. The Yale Grammatical Diversity Project documents
syntactic variation in North American English, including features such as
multiple modals in South Midland and Southern varieties
\cite{ygdp,ygdpMapbook,ygdpMultipleModals}.

For Weksa, US English should therefore be represented as a stack:

\begin{itemize}[noitemsep]
  \item \texttt{en-US.base}: broad American English defaults.
  \item \texttt{en-US.west\_coast}: a regional family containing California,
  Pacific Northwest, and urban/rural subprofiles.
  \item \texttt{en-US.east\_coast}: a family containing New York City,
  Boston, Philadelphia, Appalachian, Mid-Atlantic, and other local profiles.
  \item \texttt{en-US.deep\_south}: Southern regional profiles with separate
  gates for AAL, Appalachian, Gulf Coast, rural, urban, class, race, religion,
  and formality.
  \item \texttt{en-US.nyc.yiddish\_contact}: a New York and American Jewish
  English contact overlay, not a generic comedy seasoning.
\end{itemize}

West Coast targets need internal separation. Pacific Northwest English is an
active object of regional study, and Oregon data show that West Coast English
cannot be reduced to California alone \cite{pnwEnglish,oregonWestCoast}.
California-facing profiles should represent California Vowel Shift evidence and
its social meanings as optional phonological and indexical pressure, not as a
lexical stereotype \cite{sfCvs,villarreal2018}. A ``West Coast casual'' target
may therefore activate low formality, regional lexical choices, and media/social
stance, but it should not imply surfer speech, Valley Girl caricature, or a
single Californian identity.

Deep South targets need the same separation discipline. Multiple modals,
second-person plural forms such as \emph{y'all}, Southern vowel systems, and
regional discourse norms are real affordances, but ``Southern'' is not a class,
race, education level, or joke \cite{ygdpMultipleModals,anae}. A Deep South
stack must keep regional dialect, AAL/African American Language, Appalachian
English, Gulf Coast varieties, rurality, religion, occupation, class, and
persona as distinct axes. The machine should be able to say ``formal Black
church register in coastal Georgia'' or ``white rural Alabama mechanic talking
to a cousin'' without treating either as the whole South.

New York English also requires contact-aware modeling. New York City English
has been studied as a locally distinct variety with ethnic and neighborhood
variation, while Jewish English and Yiddish influence form their own repertoire
with lexical, prosodic, pragmatic, and identity-indexing resources
\cite{newmanNycEnglish,becker2009nyc,burdin2018,benorJewishEnglish}.
The requested ``strong New York slang with Yiddish mixed in'' should compile as
\texttt{en-US.nyc.urban} plus an American Jewish English or Yiddish-contact
overlay only when speaker, scene, community, or authoring metadata gives it
standing. The forbidden flattening is obvious and load-bearing: do not conflate
New York English, Jewish English, Yeshivish or Hasidic English, Italian
American English, Puerto Rican English, AAL, or media gangster parody.

\subsection{Scottish Targets Require Varietal Separation}

``Highland Scottish'' is not a clean linguistic target. A generation system must
separate Scots, Scottish Standard English, Highland English, and Scottish
Gaelic. The Scots Syntax Atlas documents syntactic variation across Scots
varieties with data from hundreds of speakers and many locations
\cite{scotsSyntaxAtlas}. Dictionaries of the Scots Language document Older and
Modern Scots as a historical lexicographic resource \cite{dsl}. Aitken's
Scots--English continuum remains an important conceptual anchor
\cite{aitken1984}. Highland English, by contrast, is Scottish English used in
areas shaped by Gaelic contact, especially the Highlands and Hebrides
\cite{highlandEnglish}.

For Weksa, the desired ``nigh-incomprehensible highland Scottish'' output should
not be a generic eye-dialect. It should be a selected profile stack:

\begin{itemize}[noitemsep]
  \item \texttt{en-GB.scotland.highland\_english} for Gaelic-substrate Scottish
  English;
  \item \texttt{sco.broad\_scots} or a more specific Scots variety such as
  Doric or Shetland Scots when Scots, not Highland English, is intended;
  \item an intelligibility-density control that can increase local lexical,
  syntactic, and orthographic density only when evidence supports it.
\end{itemize}

\subsection{AI Dialect Work Makes Guardrails Structural, Not Optional}

Dialect lowering must be treated as a high-risk generation task. Koenecke et
al. found racial disparities in automated speech recognition, with commercial
systems producing higher error rates for Black speakers than white speakers
\cite{koenecke2020}. Hofmann et al. show that language models can exhibit covert
racism triggered by African American English features, even when overt race
labels are absent \cite{hofmann2024}. This means Weksa must not simply ask a
model to ``make it AAVE'' or ``make it Scots.'' The state document must carry
authority, evidence, activation requirements, and forbidden uses into the
lowering path.

\subsection{Brazilian Portuguese Requires Profile Stacks, Not a Single Dialect}

Brazilian Portuguese is too internally diverse to be represented by a single
``pt-BR flavor.'' The Atlas Linguistico do Brasil (ALiB) is a national
geolinguistic anchor for regional variation and frames Brazilian Portuguese as a
multidimensional object of description \cite{alibInfo}. NURC and NURC/SP
document the cultivated urban norm in major Brazilian capitals and provide a
useful anchor for formal, media-facing, or educated urban registers
\cite{nurcsp}. Projeto SP2010 supplies a sociolinguistic corpus of Sao Paulo
speech with metadata filters for region, class, gender, age, family generation
in the city, and income \cite{sp2010}. These are the variables a lowering
system must preserve as retrieval filters rather than flattening into
``Brazilian.''

Therefore Weksa should model Brazilian targets as profile stacks:

\begin{itemize}[noitemsep]
  \item \texttt{pt-BR.base}: broad Brazilian Portuguese defaults.
  \item \texttt{pt-BR.rio.urban\_slang}: an urban Rio/carioca slang overlay,
  activated only by scene, speaker, or authoring context.
  \item \texttt{pt-BR.sao\_paulo.broadcast\_reporter}: a media-facing Sao
  Paulo register, anchored in cultivated urban norm and broadcast style rather
  than street slang.
  \item \texttt{pt-BR.minas.worker}: a Minas Gerais regional profile plus
  occupational/class persona pressure. ``Worker'' is not a dialect; it is a
  social and occupational register layered over a regional profile.
\end{itemize}

The Minas target is especially sensitive. Academic work on Mineiro speech
describes phonological, morphosyntactic, lexical, pragmatic, and historical
variation, but ``worker'' must not become shorthand for rurality, low
education, comedy, or class contempt \cite{mineiroVariation,mineiroHistoria}.
The schema must keep regional dialect, occupation, class, education, and persona
as separate axes.

\subsection{Japanese Requires Regional and Historical Axes}

Japanese target profiles need both geography and period. COJADS, the Corpus of
Japanese Dialects, is a NINJAL corpus of discourse audio from Japanese dialects
and provides a major evidence source for contemporary regional speech
\cite{cojads}. LT4All documentation describes the role of COJADS and related
language-resource work for endangered languages and dialects in Japan
\cite{lt4allJapan}. The Corpus of Historical Japanese (CHJ) provides diachronic
materials from early historical periods through later eras, while ONCOJ focuses
on Old Japanese with syntactic annotation \cite{chj,oncoj}. LAJaR adds a
typological atlas perspective across Japanese and Ryukyuan varieties
\cite{lajar2024}.

For Weksa this implies at least four independent axes:

\begin{itemize}[noitemsep]
  \item region or dialect area, such as Kansai, Tohoku, Kagoshima, Hachijo, or
  Tokyo;
  \item relationship to Ryukyuan languages and Okinawan Japanese, without
  collapsing Ryukyuan languages into Japanese dialect decoration;
  \item historical period, such as Old Japanese, Early Middle Japanese, Edo, or
  Meiji/Taisho;
  \item register and social relation, including politeness, honorific pressure,
  sentence-final particles, and stance.
\end{itemize}

Historical Japanese output should therefore use a period profile, not merely
archaic vocabulary. A profile like
\texttt{ja-JP.historical.early\_edo.commoner} would need period, region, genre,
orthographic convention, evidence corpus, and reviewer status.

\subsection{Qdrant Is Retrieval, CultCache Is Authority}

The cultural database can live in Qdrant for semantic retrieval, but Qdrant must
not own accepted language truth. Qdrant stores points with vectors and payloads;
its filtering system allows semantic search to be constrained by structured
metadata such as language, region, period, source type, and license
\cite{qdrantFiltering,qdrantPayload}. This is useful during utterance lowering:
the lowerer can retrieve candidate evidence for a requested dialect or cultural
profile.

The authority boundary should be:

\begin{lstlisting}[basicstyle=\ttfamily\small,breaklines=true]
Qdrant:
  owns: semantic retrieval over source chunks, examples, notes, and citations
  does_not_own: accepted cultural ontology profiles

CultCache:
  owns: typed cultural ontology profile documents, reviewed overlays,
        activation gates, evidence hashes, and lowering trace witnesses
\end{lstlisting}

During lowering, Weksa should assemble an evidence bundle from Qdrant, but only
a reviewed or explicitly provisional CultCache profile may authorize a dialectal
target. A Qdrant nearest neighbor is a witness candidate, not a cultural fact.

\section{Design Principles for Weksa}

\begin{enumerate}
  \item \textbf{Meaning stays upstream.} Interlingua owns referents,
  predications, relations, modality, and discourse intent.
  \item \textbf{Culture owns salience and risk.} Cultural ontology owns which
  distinctions must be inspected and which forms carry social meaning.
  \item \textbf{Dialect owns affordances, not costume.} Dialect profiles expose
  grammar, phonology, lexicon, discourse, and orthography as selectable,
  evidence-backed affordances.
  \item \textbf{Activation is gated.} A profile activates only when the speaker,
  scene, community, or authoring metadata justifies it.
  \item \textbf{Density is controlled.} Output may vary from light marking to
  dense community-internal style, but density is an explicit control.
  \item \textbf{Backtranslation is audit only.} Backtranslation helps review,
  but it must not become the source for other targets.
  \item \textbf{Forbidden flattenings are schema data.} Anti-caricature,
  anti-erasure, and anti-appropriation constraints belong in the state document,
  not in downstream apology prose.
\end{enumerate}

\section{Proposed CultCache State Shape}

The following shape should become a typed CultCache document family:
\texttt{weksa.cultural\_ontology\_profile.v0}. It can represent base languages,
dialects, sociolects, regional overlays, ritual registers, and communities of
practice.

\begin{lstlisting}[basicstyle=\ttfamily\small,breaklines=true]
schema: weksa.cultural_ontology_profile.v0
profile_id: weksa.cultural_ontology.en-US.aal.v0
kind: base_language | dialect | sociolect | regional_overlay |
      religious_overlay | community_of_practice
status: draft | reviewed | deprecated

language:
  code: en
  locale: en-US
  script: Latn
  variety_label:
  endonyms: []
  exonyms: []

parent_profiles:
  - weksa.target_cultural_ontology.en-US.v0

retrieval:
  qdrant_collection: weksa_cultural_ontology_sources
  required_payload_filters:
    language_code:
    locale:
    region:
    dialect:
    period:
    register:
    source_type:
    license:
  retrieval_policy:
    min_sources:
    require_human_review_for:
      - marginalized_or_stigmatized_dialect
      - sacred_or_ritual_register
      - historical_reconstruction

community:
  regions: []
  speech_community:
  historical_context:
  identity_authority:
  consent_notes:

activation:
  allowed_when:
    - speaker_profile_refs_profile
    - scene_context_refs_profile
    - authoring_metadata_requests_profile
  forbidden_when:
    - generic_flavor
    - mockery
    - outsider_performance_without_authority
  required_confidence: reviewed | provisional | exploratory

salience_axes:
  temporality_aspect:
    distinctions: []
    lowering_questions: []
  stance:
    distinctions: []
    lowering_questions: []
  social_relation:
    distinctions: []
    lowering_questions: []
  identity_indexing:
    distinctions: []
    lowering_questions: []
  sacred_or_ritual_context:
    distinctions: []
    lowering_questions: []

linguistic_affordances:
  grammar_features:
    - feature_id:
      status: draft | reviewed | deprecated
      domain:
        layer: grammar
        category: tense_aspect_modality | negation | agreement |
                  pronoun | word_order | discourse_particle
      description:
      formal_conditions:
        syntactic_environment: []
        semantic_inputs: []
        required_interlingua_fields: []
      realization:
        function:
        surface_strategy:
        examples: []
      activation:
        allowed_when: []
        forbidden_when: []
        requires_profile_authority:
        minimum_density:
        maximum_density:
      constraints:
        grammar_contexts: []
        register_contexts: []
        speaker_contexts: []
        incompatibilities: []
      indexicality:
        social_meanings: []
        in_group_meanings: []
        out_group_risks: []
        stereotypes_to_avoid: []
      evidence:
        source_refs: []
        qdrant_source_ids: []
        cached_source_refs: []
        source_hashes: []
        reviewer_refs: []
        confidence: low | medium | high
      uncertainty:
        open_questions: []
        known_limits: []
      trace:
        trace_label:
        must_report_when_used:
        must_report_when_rejected:
  phonology_features:
    - feature_id:
      domain:
        layer: phonology
        category: vowel_shift | consonant_realization | prosody |
                  rhythm | stress | intonation
      realization:
        ipa_hint:
        orthography_should_change:
        speech_synthesis_controls: []
      constraints:
        lexical_sets: []
        environments: []
        density_safe_for_text:
      indexicality: {}
      evidence: {}
      trace: {}
  lexicon_features:
    - feature_id:
      domain:
        layer: lexicon
        category: regional_lexeme | loanword | calque | taboo |
                  honorific | kinship | ritual_term | technical_term
      lexical_items:
        - form:
          gloss:
          sense:
          register:
          allowed_speaker_standing:
          forbidden_senses: []
          collocations: []
      activation: {}
      indexicality: {}
      evidence: {}
      trace: {}
  discourse_patterns:
    - pattern_id:
      domain:
        layer: discourse
        category: stance_marker | turn_taking | address_sequence |
                  narrative_frame | ritual_formula | humor | emphasis
      interactional_function:
      sequence_shape: []
      adjacency_constraints: []
      activation: {}
      indexicality: {}
      evidence: {}
      trace: {}
  orthography_features:
    - feature_id:
      domain:
        layer: orthography
        category: spelling_variant | script_choice | punctuation |
                  transliteration | eye_dialect_guard
      allowed_surface_forms: []
      forbidden_surface_forms: []
      renderer_notes:
      activation: {}
      evidence: {}
      trace: {}
  code_switching:
    allowed:
    constraints: []
    patterns:
      - pattern_id:
        domain:
          layer: code_switching
          category: lexical_insertion | tag_switch | alternation |
                    quotation | ritual_formula | metalinguistic
        languages: []
        matrix_language:
        embedded_language:
        switch_sites: []
        pragmatic_functions: []
        constraints:
          requires_bilingual_speaker:
          requires_community_context:
          forbidden_when: []
        evidence: {}
        trace: {}

style_controls:
  density:
    min:
    max:
    default:
  intelligibility_target:
    options: [broad_audience, local, opaque_to_outsider]
  register:
    options: []

forbidden_flattenings: []

evidence:
  corpora: []
  reference_works: []
  reviewers: []
  examples: []
  cached_source_refs: []
  qdrant_source_ids: []
  source_hashes: []
  uncertainty: []

trace_requirements:
  - activated_profile
  - axes_used
  - features_used
  - features_rejected
  - reviewer_required
\end{lstlisting}

Flat lists such as \texttt{grammar\_features: [habitual\_be]} are allowed only
as summaries or search indexes. The accepted state must be the nested feature
record: a stable feature identifier, domain, activation gates, constraints,
indexical meanings, evidence, uncertainty, and trace behavior. This lets the
lowerer ask fine-grained questions without parsing prose, such as whether a
feature requires speaker authority, whether it can be rendered in text or only
speech synthesis, which source hashes support it, and what must be reported if
the feature is rejected.

Recommended CultCache paths:

\begin{lstlisting}[basicstyle=\ttfamily\small]
.weksa/cultural-ontology/{profileId}.cc
.weksa/cultural-ontology/{baseProfileId}/overlays/{overlayId}.cc
.weksa/cultural-ontology/catalogs/{catalogId}.cc
.weksa/cultural-ontology/evidence-bundles/{bundleId}.cc
\end{lstlisting}

\section{Dialect Profile Sketches}

\subsection{AAVE/AAL}

\begin{lstlisting}[basicstyle=\ttfamily\small,breaklines=true]
schema: weksa.cultural_ontology_profile.v0
profile_id: weksa.cultural_ontology.en-US.aal.v0
kind: dialect
status: draft
parent_profiles:
  - weksa.target_cultural_ontology.en-US.v0
activation:
  allowed_when:
    - speaker_profile_refs_aal
    - scene_context_refs_aal_community
    - authoring_metadata_requests_reviewed_aal
  forbidden_when:
    - generic_black_coding
    - comedy_accent
    - criminality_or_low_intelligence_indexing
salience_axes:
  temporality_aspect:
    distinctions:
      - habituality
      - remote_past_or_remote_completion
      - completive_result_state
  identity_indexing:
    distinctions:
      - in_group_style
      - out_group_performance_risk
      - bidialectal_context_switching
linguistic_affordances:
  grammar_features:
    - feature_id: aal_invariant_habitual_be
      status: draft_guarded
      domain:
        layer: grammar
        category: tense_aspect_modality
        subcategory: habitual_aspect
      description: Marks habitual or characteristic occurrence, not present
        progressive meaning.
      formal_conditions:
        semantic_inputs:
          - habituality
        required_interlingua_fields:
          - event_iterativity_or_characteristic_state
      realization:
        function: express_habitual_aspect
        surface_strategy: invariant_be
        examples: []
      activation:
        allowed_when:
          - speaker_profile_refs_aal
          - scene_context_refs_aal_community
        forbidden_when:
          - generic_black_coding
          - outsider_performance_without_authority
        requires_profile_authority: true
        minimum_density: reviewed_light
        maximum_density: reviewed_dense
      constraints:
        grammar_contexts:
          - only_when_habitual_aspect_is_intended
        incompatibilities:
          - present_progressive_only_reading
      indexicality:
        social_meanings:
          - aal_systematic_aspect
        out_group_risks:
          - treated_as_broken_standard_english
        stereotypes_to_avoid:
          - ignorance
          - comedy
      evidence:
        source_refs:
          - Green2002
          - Rickford1999
          - YaleInvariantBe
          - CORAAL
        confidence: high
      trace:
        trace_label: aal_habitual_be
        must_report_when_used: true
        must_report_when_rejected: true
    - feature_id: aal_stressed_bin
      status: draft_guarded
      domain:
        layer: grammar
        category: tense_aspect_modality
        subcategory: remote_past_or_remote_state
      activation:
        allowed_when:
          - speaker_profile_refs_aal
        forbidden_when:
          - generic_flavor
        requires_profile_authority: true
      evidence:
        source_refs:
          - Green2002
          - Rickford1999
        confidence: medium
      trace:
        trace_label: aal_stressed_bin
        must_report_when_used: true
        must_report_when_rejected: true
    - feature_id: aal_completive_done
      status: draft_guarded
      domain:
        layer: grammar
        category: tense_aspect_modality
        subcategory: completive
      activation:
        allowed_when:
          - speaker_profile_refs_aal
        forbidden_when:
          - generic_flavor
        requires_profile_authority: true
      evidence:
        source_refs:
          - Green2002
          - Rickford1999
        confidence: medium
      trace:
        trace_label: aal_completive_done
        must_report_when_used: true
        must_report_when_rejected: true
    - feature_id: aal_negative_concord
      status: draft_guarded
      domain:
        layer: grammar
        category: negation
        subcategory: negative_concord
      activation:
        allowed_when:
          - speaker_profile_refs_aal
        forbidden_when:
          - generic_flavor
        requires_profile_authority: true
      evidence:
        source_refs:
          - Green2002
          - Rickford1999
        confidence: medium
      trace:
        trace_label: aal_negative_concord
        must_report_when_used: true
        must_report_when_rejected: true
  phonology_features: []
  lexicon_features: []
  discourse_patterns: []
  orthography_features:
    - feature_id: aal_eye_dialect_guard
      status: draft_guarded
      domain:
        layer: orthography
        category: eye_dialect_guard
      allowed_surface_forms: []
      forbidden_surface_forms:
        - misspelling_as_dialect
        - random_apostrophe_deletion
      renderer_notes: Text rendering must not fake AAL phonology through
        stigmatizing spelling unless an orthographic practice is separately
        evidenced and reviewed.
      trace:
        trace_label: aal_eye_dialect_guard
        must_report_when_used: false
        must_report_when_rejected: true
  code_switching:
    allowed: true
    constraints:
      - Requires speaker/community context and evidence-backed pattern.
    patterns: []
forbidden_flattenings:
  - Do not treat AAL as broken Standard English.
  - Do not use AAL to imply criminality, ignorance, comedy, or informality.
  - Do not generate dense AAL without speaker/context authority.
\end{lstlisting}

\subsection{US English Regional and Contact Overlays}

\begin{lstlisting}[basicstyle=\ttfamily\small,breaklines=true]
schema: weksa.cultural_ontology_profile.v0
profile_id: weksa.cultural_ontology.en-US.west_coast.california_casual.v0
kind: regional_register_overlay
status: draft_guarded
parent_profiles:
  - weksa.target_cultural_ontology.en-US.v0
retrieval:
  qdrant_collection: weksa_cultural_ontology_sources
  required_payload_filters:
    language_code: en
    locale: en-US
    region: [California, West Coast]
    source_type: [dialect_atlas, sociolinguistic_study, reviewed_example]
activation:
  allowed_when:
    - speaker_or_scene_refs_california_or_west_coast
    - authoring_metadata_requests_california_casual_register
  forbidden_when:
    - surfer_or_valley_girl_caricature
    - generic_american_youth_flavor
style_controls:
  density:
    min: light
    max: reviewed_moderate
  register:
    options: [casual, media_casual, professional_informal]
forbidden_flattenings:
  - Do not treat California as the whole West Coast.
  - Do not treat vowel-shift evidence as mandatory spelling distortion.
  - Do not collapse region, age, ethnicity, and class into one style.
\end{lstlisting}

\begin{lstlisting}[basicstyle=\ttfamily\small,breaklines=true]
schema: weksa.cultural_ontology_profile.v0
profile_id: weksa.cultural_ontology.en-US.deep_south.v0
kind: regional_profile_stack_template
status: draft_guarded
parent_profiles:
  - weksa.target_cultural_ontology.en-US.v0
stack_axes:
  region: Deep South
  race_or_ethnolinguistic_tradition: explicit_or_unknown
  rural_urban: explicit_or_unknown
  class_or_occupation: explicit_or_unknown
  religion_or_church_register: explicit_or_unknown
  formality: explicit_or_unknown
activation:
  allowed_when:
    - speaker_or_scene_refs_deep_south
    - authoring_metadata_requests_specific_southern_region
  forbidden_when:
    - generic_redneck_caricature
    - low_intelligence_indexing
    - race_or_class_inferred_from_region_only
salience_axes:
  modality:
    distinctions:
      - multiple_modals_when_evidence_supports
  address:
    distinctions:
      - second_person_plural
      - politeness_and_kinship_terms
forbidden_flattenings:
  - Southern is not AAL, and AAL is not Southern by default.
  - Southern is not one accent, class, politics, religion, or education level.
  - Do not use eye spelling as a substitute for grammar or corpus evidence.
\end{lstlisting}

\begin{lstlisting}[basicstyle=\ttfamily\small,breaklines=true]
schema: weksa.cultural_ontology_profile.v0
profile_id: weksa.cultural_ontology.en-US.east_coast.v0
kind: regional_family_profile
status: draft
parent_profiles:
  - weksa.target_cultural_ontology.en-US.v0
child_profiles:
  - weksa.cultural_ontology.en-US.nyc.urban.v0
  - weksa.cultural_ontology.en-US.boston.v0
  - weksa.cultural_ontology.en-US.philadelphia.v0
  - weksa.cultural_ontology.en-US.appalachian.v0
activation:
  allowed_when:
    - specific_east_coast_region_selected
  forbidden_when:
    - generic_east_coast_attitude
forbidden_flattenings:
  - East Coast is a routing family, not a dialect.
  - Do not let New York stand in for all eastern US speech.
\end{lstlisting}

\begin{lstlisting}[basicstyle=\ttfamily\small,breaklines=true]
schema: weksa.cultural_ontology_profile.v0
profile_id: weksa.cultural_ontology.en-US.nyc.yiddish_contact_slang.v0
kind: contact_language_overlay
status: draft_guarded
parent_profiles:
  - weksa.cultural_ontology.en-US.nyc.urban.v0
contact_languages:
  - yi
speech_traditions:
  - American Jewish English
  - Yiddish contact in New York English
activation:
  allowed_when:
    - speaker_profile_refs_jewish_english_or_yiddish_contact
    - scene_context_refs_nyc_jewish_community_or_yiddish_contact
    - authoring_metadata_requests_reviewed_contact_overlay
  forbidden_when:
    - generic_new_york_comedy
    - yiddish_as_random_spice
    - non_jewish_speaker_performance_without_context
style_controls:
  density:
    min: light_lexical_color
    max: reviewed_dense_contact_register
  register:
    options: [casual, comic_in_group, family_argument, public_nyc_slang]
forbidden_flattenings:
  - Do not conflate New York English with Jewish English.
  - Do not conflate Jewish English, Yeshivish, Hasidic English, or Yiddish.
  - Do not use Yiddish-origin forms to mark greed, neurosis, or cowardice.
\end{lstlisting}

\subsection{Highland English / Gaelic-Substrate Scottish English}

\begin{lstlisting}[basicstyle=\ttfamily\small,breaklines=true]
schema: weksa.cultural_ontology_profile.v0
profile_id: weksa.cultural_ontology.en-GB.scotland.highland_english.v0
kind: regional_overlay
status: draft
parent_profiles:
  - weksa.target_cultural_ontology.en-GB.v0
contact_substrate:
  - Scottish Gaelic
activation:
  allowed_when:
    - speaker_or_scene_refs_highlands_or_hebrides
    - speaker_profile_refs_gaelic_substrate_english
  forbidden_when:
    - generic_scottish_flavor
    - pretending_scots_and_gaelic_are_the_same
style_controls:
  density:
    min: light
    max: heavy_reviewed
  intelligibility_target:
    options: [broad_audience, local, opaque_to_outsider]
forbidden_flattenings:
  - Do not collapse Highland English into Scots.
  - Do not collapse Scots into Scottish Gaelic.
  - Do not use eye-dialect gibberish as authenticity.
\end{lstlisting}

\subsection{Broad Scots / Doric / Shetland as Separate Profiles}

\begin{lstlisting}[basicstyle=\ttfamily\small,breaklines=true]
schema: weksa.cultural_ontology_profile.v0
profile_id: weksa.cultural_ontology.sco.broad_scots.v0
kind: dialect_or_language_profile
status: draft
language:
  code: sco
  locale: Scotland
parent_profiles:
  - weksa.cultural_ontology.en-GB.scotland.continuum.v0
activation:
  allowed_when:
    - scene_or_speaker_refs_scots
    - authoring_metadata_selects_scots_variety
  forbidden_when:
    - generic_highland_english
    - generic_accent_joke
style_controls:
  density:
    min: light_scots_lexis
    max: broad_reviewed_scots
  intelligibility_target:
    options: [broad_audience, local, opaque_to_outsider]
evidence:
  corpora:
    - ScotsSyntaxAtlas
  reference_works:
    - DictionariesOfTheScotsLanguage
\end{lstlisting}

\subsection{Brazilian Regional Targets}

\begin{lstlisting}[basicstyle=\ttfamily\small,breaklines=true]
schema: weksa.cultural_ontology_profile.v0
profile_id: weksa.cultural_ontology.pt-BR.rio.urban_slang.v0
kind: regional_overlay
status: draft_guarded
parent_profiles:
  - weksa.target_cultural_ontology.pt-BR.v0
activation:
  allowed_when:
    - scene_or_speaker_refs_rio_urban_context
    - authoring_metadata_requests_carioca_urban_register
  forbidden_when:
    - generic_brazilian_flavor
    - comedy_crime_or_hypersexualized_indexing
retrieval:
  qdrant_collection: weksa_cultural_ontology_sources
  required_payload_filters:
    language_code: pt
    locale: pt-BR
    region: Rio de Janeiro
    source_type: [corpus, sociolinguistic_study, reviewed_example]
style_controls:
  density:
    min: light
    max: reviewed_dense
forbidden_flattenings:
  - Do not treat Rio slang as all Brazilian Portuguese.
  - Do not use slang density to imply criminality or unseriousness.
\end{lstlisting}

\begin{lstlisting}[basicstyle=\ttfamily\small,breaklines=true]
schema: weksa.cultural_ontology_profile.v0
profile_id: weksa.cultural_ontology.pt-BR.sao_paulo.broadcast_reporter.v0
kind: institutional_register_overlay
status: draft
parent_profiles:
  - weksa.target_cultural_ontology.pt-BR.v0
evidence:
  corpora:
    - NURC_SP
    - Projeto_SP2010
activation:
  allowed_when:
    - speaker_role_is_broadcast_reporter
    - scene_medium_is_news_or_public_institutional_speech
style_controls:
  density:
    default: low_regional_marking
  register:
    options: [broadcast_neutral, formal_live_report, conversational_anchor]
forbidden_flattenings:
  - Do not confuse Sao Paulo media register with all Sao Paulo speech.
  - Do not erase intentional informalization when the scene calls for it.
\end{lstlisting}

\begin{lstlisting}[basicstyle=\ttfamily\small,breaklines=true]
schema: weksa.cultural_ontology_profile.v0
profile_id: weksa.cultural_ontology.pt-BR.minas.worker.v0
kind: profile_stack_template
status: draft_guarded
parent_profiles:
  - weksa.target_cultural_ontology.pt-BR.v0
stack_axes:
  region: Minas Gerais
  occupation_or_class_register: worker
  education: explicit_or_unknown
  rural_urban: explicit_or_unknown
activation:
  allowed_when:
    - speaker_profile_refs_minas_region
    - speaker_profile_refs_worker_or_occupation
  forbidden_when:
    - class_caricature
    - rural_bumpkin_default
    - low_intelligence_indexing
forbidden_flattenings:
  - Worker is not a dialect.
  - Mineiro is not one uniform speech form.
  - Do not use eye spelling as a substitute for phonology or grammar evidence.
\end{lstlisting}

\subsection{Japanese Regional and Historical Targets}

\begin{lstlisting}[basicstyle=\ttfamily\small,breaklines=true]
schema: weksa.cultural_ontology_profile.v0
profile_id: weksa.cultural_ontology.ja-JP.kansai.casual.v0
kind: regional_overlay
status: draft
parent_profiles:
  - weksa.target_cultural_ontology.ja-JP.v0
retrieval:
  qdrant_collection: weksa_cultural_ontology_sources
  required_payload_filters:
    language_code: ja
    region: Kansai
    period: contemporary
    source_type: [corpus, dialect_reference, reviewed_example]
salience_axes:
  sentence_final_force:
  politeness:
  stance:
  humor_or_informality_indexing:
\end{lstlisting}

\begin{lstlisting}[basicstyle=\ttfamily\small,breaklines=true]
schema: weksa.cultural_ontology_profile.v0
profile_id: weksa.cultural_ontology.ja-JP.historical.early_edo.v0
kind: historical_period_profile
status: draft_guarded
parent_profiles:
  - weksa.target_cultural_ontology.ja-JP.v0
retrieval:
  qdrant_collection: weksa_cultural_ontology_sources
  required_payload_filters:
    language_code: ja
    period: Edo
    source_type: [historical_corpus, scholarly_reference]
activation:
  allowed_when:
    - scene_period_refs_edo
    - authoring_metadata_requests_historical_japanese
  forbidden_when:
    - generic_archaic_flavor
    - anime_period_speech_without_source
evidence:
  corpora:
    - Corpus_of_Historical_Japanese
    - ONCOJ_when_old_japanese_not_edo
forbidden_flattenings:
  - Do not use modern dialect stereotypes as historical Japanese.
  - Do not mix periods unless the source intentionally codes anachronism.
\end{lstlisting}

\section{Implementation Consequences}

Weksa should not implement dialect lowering as a prompt modifier. The lowering
path should compile a profile stack:

\begin{lstlisting}[basicstyle=\ttfamily\small]
interlingua packet
  -> projected speaker context
  -> cultural ontology profile stack
  -> dialect/register realization plan
  -> target utterance
  -> trace naming profile axes and features used
\end{lstlisting}

The cultural ontology profile should be inspectable as CultCache state before
runtime lowering. A target utterance trace should name:

\begin{itemize}[noitemsep]
  \item which profile stack activated;
  \item which salience axes mattered;
  \item which grammatical or lexical affordances were used;
  \item which affordances were rejected;
  \item whether human review is required.
\end{itemize}

\section{Conclusion}

The contemporary Sapir--Whorf lesson for Weksa is not that culture determines
thought. It is that culture and speech-community practice shape what must be
noticed, what can be said naturally, what indexes identity, and what becomes
dangerous when imitated without authority. A useful cultural ontology schema
therefore represents salience, indexicality, ideology, affordance, activation,
evidence, and forbidden flattening. That schema is the missing bridge between
language-neutral interlingua and plausible, reviewable dialectal utterance.

\section*{Citation Ledger}

\begin{longtable}{>{\raggedright\arraybackslash}p{0.28\linewidth}
                  >{\raggedright\arraybackslash}p{0.66\linewidth}}
\textbf{Source} & \textbf{Use in Weksa synthesis} \\
\hline
Slobin 1996 & Supports the ``thinking for speaking'' mechanism: language
guides attention during utterance formulation, not total cognition. \\
Lucy 1992, 1997 & Provides the careful reformulation of linguistic relativity
as empirical comparison of language categories and interpretation. \\
Gumperz and Levinson 1996 & Frames relativity as language, culture, and
classification, with conversational inference as a key cultural site. \\
Wolff and Holmes 2011 & Contemporary cognitive review: relativity is contested
but supported in bounded mechanisms and domains. \\
Silverstein 2003 & Supplies indexical order: forms presuppose and create
social context at multiple levels. \\
Eckert 2012 & Treats variation as social meaning, supporting dialect profiles
as more than demographic labels. \\
Bucholtz and Hall 2005 & Identity is interactionally produced and indexed by
linguistic practice; supports activation gates and speaker authority. \\
Irvine and Gal 2000 & Language ideology mechanisms warn against erasure,
iconization, and flattening. \\
Green 2002; Rickford 1999, 2016 & Establish AAVE/AAL as systematic grammar and
socially meaningful language practice. \\
CORAAL, ORAAL, Yale Grammatical Diversity Project & Provide corpus and pattern
anchors for AAL features and reviewed examples. \\
DARE, ANAE, YGDP & Anchor US English as regional lexical, phonological, and
syntactic variation rather than a single American default. \\
Pacific Northwest English Study; Becker et al.; Hall-Lew et al.; Villarreal &
Anchor West Coast internal variation, including Oregon, Pacific Northwest, San
Francisco, and California Vowel Shift social meanings. \\
Newman; Becker; Burdin; Benor and Cohen & Anchor New York City English,
ethnicity in NYC speech, Jewish English prosody, and Yiddish/American Jewish
English contact as separate but intersecting resources. \\
Scots Syntax Atlas, DSL, Aitken 1984 & Provide dialectological and
lexicographic anchors for Scots and the Scots--English continuum. \\
Highland English reference & Warns that Highland English is Gaelic-contact
Scottish English, not simply Scots. \\
ALiB, NURC/SP, Projeto SP2010 & Anchor Brazilian Portuguese variation,
cultivated urban norms, and Sao Paulo sociolinguistic metadata. \\
Mineiro studies & Warn that Minas Gerais speech must be represented as
regional variation plus social/persona axes, not as class caricature. \\
COJADS, CHJ, ONCOJ, LAJaR & Anchor contemporary Japanese dialect evidence,
historical Japanese evidence, Old Japanese annotation, and Japanese/Ryukyuan
typological coverage. \\
Qdrant documentation & Supports the retrieval architecture: vectors plus
payload filters for semantic evidence search. \\
Koenecke et al. 2020; Hofmann et al. 2024 & Show automated systems fail and
discriminate around dialects; justify structural guardrails. \\
\end{longtable}

\begin{thebibliography}{99}

\bibitem{alibInfo}
Projeto Atlas Linguistico do Brasil (ALiB).
``The Linguistic Atlas of Brazil.''
Universidade Federal da Bahia.
\url{https://alib.ufba.br/sites/alib.ufba.br/files/informacoes_em_ingles.pdf}.

\bibitem{aitken1984}
Aitken, A. J. 1984.
``Scots and English in Scotland.''
In Peter Trudgill (ed.), \emph{Language in the British Isles}.
Cambridge: Cambridge University Press.
Available via the Scots Language Centre:
\url{https://media.scotslanguage.com/library/document/aitken/Scots_and_English_in_Scotland.pdf}.

\bibitem{anae}
Labov, William, Sharon Ash, and Charles Boberg. 2006.
\emph{The Atlas of North American English: Phonetics, Phonology and Sound
Change}.
Berlin: Mouton de Gruyter.
\url{https://doi.org/10.1515/9783110206838}.

\bibitem{becker2009nyc}
Becker, Kara. 2009.
``The Sociolinguistics of Ethnicity in New York City.''
\emph{Language and Linguistics Compass} 3(3):751--766.
\url{https://doi.org/10.1111/j.1749-818X.2009.00138.x}.

\bibitem{benorJewishEnglish}
Benor, Sarah Bunin, and Steven M. Cohen. 2009.
``Talking Jewish: The `ethnic English' of American Jews.''
In Lederhendler (ed.), \emph{Ethnicity and Beyond: Theories and Dilemmas of
Jewish Group Demarcation}, 62--78.

\bibitem{burdin2018}
Burdin, Rachel Steindel. 2018.
``The Perception of Macro-rhythm in Jewish English Intonation.''
\emph{Penn Working Papers in Linguistics} 24(2).
\url{https://scholars.unh.edu/faculty_pubs/1459/}.

\bibitem{bucholtzhall2005}
Bucholtz, Mary, and Kira Hall. 2005.
``Identity and interaction: A sociocultural linguistic approach.''
\emph{Discourse Studies} 7(4--5):585--614.
\url{https://doi.org/10.1177/1461445605054407}.

\bibitem{coraal}
Kendall, Tyler, and Charlie Farrington.
\emph{Corpus of Regional African American Language (CORAAL)}.
Online Resources for African American Language.
\url{https://oraal.github.io/coraal}.

\bibitem{chj}
National Institute for Japanese Language and Linguistics (NINJAL).
\emph{Corpus of Historical Japanese (CHJ): Overview}.
\url{https://clrd.ninjal.ac.jp/chj/overview-en.html}.

\bibitem{cojads}
National Institute for Japanese Language and Linguistics (NINJAL).
\emph{Corpus of Japanese Dialects (COJADS)}.
\url{https://www2.ninjal.ac.jp/cojads/index.html}.

\bibitem{dare}
Dictionary of American Regional English.
``What is DARE?''
University of Wisconsin--Madison.
\url{https://dare.wisc.edu/about/what-is-dare/}.

\bibitem{dareReport}
von Schneidemesser, Luanne. 2006.
``The Dictionary of American Regional English: A Report on its Present State
and Plans.''
In \emph{Proceedings of EURALEX 2006}.
\url{https://euralex.org/elx_proceedings/Euralex2006/060_2006_V1_Luanne%20von%20SCHNEIDEMESSER_The%20Dictionary%20of%20American%20Regional%20English_a%20Report%20on%20its.pdf}.

\bibitem{dsl}
Dictionaries of the Scots Language.
\emph{DSL Online}.
\url{https://dsl.ac.uk/our-publications/what-is-dsl/}.

\bibitem{eckert2012}
Eckert, Penelope. 2012.
``Three waves of variation study: The emergence of meaning in the study of
variation.''
\emph{Annual Review of Anthropology} 41:87--100.
\url{https://doi.org/10.1146/annurev-anthro-092611-145828}.

\bibitem{oregonWestCoast}
Becker, Kara, et al.
``Variation in West Coast English: The Case of Oregon.''
\emph{Publication of the American Dialect Society}.
\url{https://www.reed.edu/linguistics/becker/Becker_et_al_PADS_inpress.pdf}.

\bibitem{green2002}
Green, Lisa J. 2002.
\emph{African American English: A Linguistic Introduction}.
Cambridge: Cambridge University Press.
\url{https://www.cambridge.org/core/books/african-american-english/4CE71E97BA60212C94A342CC7F912250}.

\bibitem{gumperzlevinson1996}
Gumperz, John J., and Stephen C. Levinson, eds. 1996.
\emph{Rethinking Linguistic Relativity}.
Cambridge: Cambridge University Press.
\url{https://books.google.com/books/about/Rethinking_Linguistic_Relativity.html?id=o7U7PgAACAAJ}.

\bibitem{highlandEnglish}
McArthur, Tom, ed.
``Highland English.''
\emph{Concise Oxford Companion to the English Language}, entry reproduced at
Encyclopedia.com.
\url{https://www.encyclopedia.com/humanities/encyclopedias-almanacs-transcripts-and-maps/highland-english}.

\bibitem{hofmann2024}
Hofmann, Valentin, Pratyusha Ria Kalluri, Dan Jurafsky, and Sharese King. 2024.
``AI generates covertly racist decisions about people based on their dialect.''
\emph{Nature} 633:147--154.
\url{https://doi.org/10.1038/s41586-024-07856-5}.

\bibitem{irvinegal2000}
Irvine, Judith T., and Susan Gal. 2000.
``Language ideology and linguistic differentiation.''
In Paul V. Kroskrity (ed.), \emph{Regimes of Language: Ideologies, Polities,
and Identities}, 35--83.
Santa Fe: School of American Research Press.

\bibitem{koenecke2020}
Koenecke, Allison, Andrew Nam, Emily Lake, Joe Nudell, Minnie Quartey, Zion
Mengesha, Connor Toups, John R. Rickford, Dan Jurafsky, and Sharad Goel. 2020.
``Racial disparities in automated speech recognition.''
\emph{Proceedings of the National Academy of Sciences} 117(14):7684--7689.
\url{https://doi.org/10.1073/pnas.1915768117}.

\bibitem{lucy1992}
Lucy, John A. 1992.
\emph{Language Diversity and Thought: A Reformulation of the Linguistic
Relativity Hypothesis}.
Cambridge: Cambridge University Press.

\bibitem{lucy1997}
Lucy, John A. 1997.
``Linguistic relativity.''
\emph{Annual Review of Anthropology} 26:291--312.
\url{https://www.linguisticanthropology.org/wp-content/uploads/2010/08/1997-Lucy-LinguisticRelativity-s.pdf}.

\bibitem{lajar2024}
Kato, Kanji, So Miyagawa, and Natsuko Nakagawa. 2024.
``Language Atlas of Japanese and Ryukyuan (LAJaR): A Linguistic Typology
Database for Endangered Japonic Languages.''
In \emph{Proceedings of the 6th Workshop on Research in Computational
Linguistic Typology and Multilingual NLP}, 55--57.
Association for Computational Linguistics.
\url{https://doi.org/10.18653/v1/2024.sigtyp-1.7}.

\bibitem{lt4allJapan}
National Institute for Japanese Language and Linguistics (NINJAL). 2019.
``Language Resources for Endangered Languages and Dialects in Japan.''
\emph{Proceedings of Language Technologies for All (LT4All)}, 371--374.
\url{https://lt4all.elra.info/proceedings/lt4all2019/pdf/2019.lt4all-1.93.pdf}.

\bibitem{newmanNycEnglish}
Newman, Michael. 2020.
``New York City English.''
In \emph{The Handbook of World Englishes}, 2nd ed.
Wiley.
\url{https://doi.org/10.1002/9781119518297.eowe00153}.

\bibitem{mineiroHistoria}
Ribeiro, Silvana Soares Costa. 2016.
``A relacao entre dialetologia e historia: reflexoes teorico-metodologicas
para o estudo do portugues usado em Minas Gerais.''
\emph{Gragoata}.
\url{https://periodicos.uff.br/gragoata/article/view/33374}.

\bibitem{mineiroVariation}
Oliveira, Gisele et al.
``A variacao linguistica do `mineires': tracos fonetico-fonologicos,
morfossintaticos, lexicais e pragmaticos em perspectiva sociolinguistica.''
\emph{Sumarios.org record}.
\url{https://sumarios.org/artigo/varia%C3%A7%C3%A3o-lingu%C3%ADstica-do-%E2%80%9Cmineir%C3%AAs%E2%80%9D-tra%C3%A7os-fon%C3%A9tico-fonol%C3%B3gicos-morfossint%C3%A1ticos-lexicais-e}.

\bibitem{nurcsp}
Pellegrino, Alexandre, Aline B. Mendes, and Marcelo Finger. 2022.
``Bringing NURC/SP to Digital Life: the Role of Open-source Automatic Speech
Recognition Models.''
\url{https://arxiv.org/abs/2210.07852}.

\bibitem{oncoj}
University of Oxford and NINJAL.
\emph{Oxford-NINJAL Corpus of Old Japanese (ONCOJ)}.
\url{https://oncoj.ninjal.ac.jp/}.

\bibitem{oraalPatterns}
Online Resources for African American Language.
``African American Language Linguistic Patterns.''
\url{https://oraal.github.io/aal-linguistic-patterns}.

\bibitem{pnwEnglish}
Pacific Northwest English Study.
``Project Description.''
University of Washington.
\url{https://zeos.ling.washington.edu/PNWEnglish/about/projectDescription.php}.

\bibitem{qdrantFiltering}
Qdrant.
``Filtering.''
\url{https://qdrant.tech/documentation/search/filtering/}.

\bibitem{qdrantPayload}
Qdrant.
``A Complete Guide to Filtering in Vector Search.''
\url{https://qdrant.tech/articles/vector-search-filtering/}.

\bibitem{rickford1999}
Rickford, John R. 1999.
\emph{African American Vernacular English: Features, Evolution, Educational
Implications}.
Oxford: Blackwell.
\url{https://www.wiley-vch.de/de/fachgebiete/geistes-und-sozialwissenschaften/african-american-vernacular-english-978-0-631-21245-4}.

\bibitem{rickford2016}
Rickford, John R. 2016.
``Labov's contributions to the study of African American Vernacular English:
Pursuing linguistic and social equity.''
\emph{Journal of Sociolinguistics} 20(4):561--580.
\url{https://doi.org/10.1111/josl.12198}.

\bibitem{scotsSyntaxAtlas}
Scots Syntax Atlas.
``Research context'' and ``About the data.''
University of Glasgow project.
\url{https://scotssyntaxatlas.ac.uk/research-context/};
\url{https://scotssyntaxatlas.ac.uk/about/about-the-data/}.

\bibitem{sp2010}
Projeto SP2010.
``Corpus.''
Universidade de Sao Paulo.
\url{https://projetosp2010.fflch.usp.br/corpus}.

\bibitem{silverstein2003}
Silverstein, Michael. 2003.
``Indexical order and the dialectics of sociolinguistic life.''
\emph{Language \& Communication} 23(3--4):193--229.
\url{https://doi.org/10.1016/S0271-5309(03)00013-2}.

\bibitem{sfCvs}
Hall-Lew, Lauren, Amanda Cardoso, Yova Kemenchedjieva, Kieran Wilson, Ruaridh
Purse, and Julie Saigusa. 2015.
``San Francisco English and the California Vowel Shift.''
Proceedings of ICPhS 2015.
\url{https://www.pure.ed.ac.uk/ws/portalfiles/portal/21806704/Hall_Lew_etal_2015_ICPhS_paper.pdf}.

\bibitem{slobin1996}
Slobin, Dan I. 1996.
``From `thought and language' to `thinking for speaking'.''
In John J. Gumperz and Stephen C. Levinson (eds.),
\emph{Rethinking Linguistic Relativity}, 70--96.
Cambridge: Cambridge University Press.

\bibitem{wolffholmes2011}
Wolff, Phillip, and Kevin J. Holmes. 2011.
``Linguistic relativity.''
\emph{Wiley Interdisciplinary Reviews: Cognitive Science} 2(3):253--265.
\url{https://doi.org/10.1002/wcs.104}.

\bibitem{villarreal2018}
Villarreal, Dan. 2018.
``The construction of social meaning: A matched-guise investigation of the
California Vowel Shift.''
\emph{Journal of English Linguistics}.
\url{https://doi.org/10.1177/0075424217753520}.

\bibitem{ygdp}
Yale Grammatical Diversity Project.
``Project Description.''
\url{https://ygdp.yale.edu/project-description}.

\bibitem{yaleInvariantBe}
Yale Grammatical Diversity Project.
``Invariant be.''
\url{https://ygdp.yale.edu/phenomena/invariant-be}.

\bibitem{ygdpMapbook}
Wood, Jim, Kaija Gahm, Raffaella Zanuttini, Laurence R. Horn, and Jason
Kandybowicz. 2019.
\emph{The Yale Grammatical Diversity Project: A Mapbook of Syntactic Variation
in American English}.
\url{https://elischolar.library.yale.edu/ygdp/7/}.

\bibitem{ygdpMultipleModals}
Yale Grammatical Diversity Project.
``Multiple modals.''
\url{https://ygdp.yale.edu/phenomena/multiple-modals}.

\end{thebibliography}

\end{document}
